
\medskip
\noindent\SGAInlineRefTarget{sga3:vii:paragraph:1-2-4}\textbf{1.2.4.}\footnote{Editors' note: In this paragraph the
order has been changed: \(\widehat S_k(E)\) is defined first, and it is
then stated that \(\mathbf V^f_k(E)\) represents
\(\mathbf V^f_k(E)\).}
\phantomsection
\label{sga3:VIIB:1.2.4}
For every pseudocompact \(k\)-module \(E\), let
\(\widehat S_k(E)\) denote the \emph{completed symmetric algebra} of
\(E\), defined as follows. Let \(T_k(E)\) be the direct sum of the
pseudocompact \(k\)-modules
\[
  \widehat{\bigotimes}_{k}^{\,n}E
  =
  E\widehat{\otimes}_{k}\cdots\widehat{\otimes}_{k}E
  \qquad
  \left(
    n\geq0;\quad
    \widehat{\bigotimes}_{k}^{\,0}E=k
  \right).
\]
Give \(T_k(E)\) its usual graded \(k\)-algebra multiplication. Let
\(S_k(E)\) be the quotient of \(T_k(E)\) by the homogeneous ideal
whose \(n\)-th component is the closed \(k\)-submodule of
\(\widehat{\bigotimes}_{k}^{\,n}E\) generated by the elements
\[
\begin{aligned}
&x_1\widehat{\otimes}\cdots\widehat{\otimes}x_i
 \widehat{\otimes}x_{i+1}\widehat{\otimes}\cdots
 \widehat{\otimes}x_n
\\[-0.1em]
&\qquad{}-
x_1\widehat{\otimes}\cdots\widehat{\otimes}x_{i+1}
 \widehat{\otimes}x_i\widehat{\otimes}\cdots
 \widehat{\otimes}x_n.
\end{aligned}
\]
If this quotient is denoted by \(S_k^n(E)\), then \(S_k(E)\) is a
graded \(k\)-algebra whose \(n\)-th component is \(S_k^n(E)\).

Endow \(S_k(E)\) with the linear topology defined by the set
\(\Upsilon\) of ideals \(\mathfrak l\) such that
\(S_k(E)/\mathfrak l\) is a \(k\)-module of finite length and
\[
  S_k^n(E)\cap\mathfrak l
\]
is an open submodule of \(S_k^n(E)\) for every \(n\). The profinite
algebra \(\widehat S_k(E)\) is the separated completion of \(S_k(E)\)
for this topology.\footnote{Editors' note: For example, let \(k\) be
a field, let \(E\) be a pseudocompact \(k\)-vector space, and put
\(V=\operatorname{Hom}_c(E,k)\); then \(E\simeq V^*\)
\textup{(cf.\ \SGARefLink{sga3:VIIB:0.2.2}{0.2.2})}. For every finite-dimensional subspace \(W\) of
\(V\), let \(\mathcal F(W)\) be the set of closed points of the
\(k\)-scheme
\(\mathbf V(W)=\operatorname{Spec}S_k(W^*)\), and let
\(\mathcal F(V)=\varinjlim_W\mathcal F(W)\). Then
\(\widehat S_k(E)\) is the product, over \(x\in\mathcal F(V)\), of
the pseudocompact \(k\)-algebras
\[
  \widehat S_k(E)_x
  =
  \varprojlim_W\widehat{\mathcal O}_{W,x},
\]
where \(W\) ranges over the finite-dimensional subspaces of \(V\)
such that \(x\in\mathcal F(W)\), and
\(\widehat{\mathcal O}_{W,x}\) is the separated completion of the
local ring \(\mathcal O_{\mathbf V(W),x}\) for the topology defined
by the ideals of finite codimension \textup{(which here coincides
with the \(\mathfrak m\)-adic topology)}. If \(k\) is algebraically
closed and \((v_i)_{i\in I}\) is a basis of \(V\), so that \(E\) has
a pseudobasis \((e_i)_{i\in I}\), then every local component
\(\widehat S_k(E)_x\) is isomorphic to the formal power-series ring
\(k[[e_i,\ i\in I]]\), endowed with the topology defined by the
ideals of finite codimension.}

We write
\[
  \mathbf V^f_k(E)
  =
  \operatorname{Spf}\bigl(\widehat S_k(E)\bigr)
\]
for the corresponding formal \(k\)-variety. It represents the
\(k\)-functor \(\mathbf V^f_k(E)\); that is, for every object \(C\)
of \(\operatorname{Alf}/k\), there is a canonical isomorphism
\[
\begin{aligned}
  \mathbf V^f_k(E)(C)
  &=
  \operatorname{Hom}_c(E,C|_k)
  \xrightarrow{\ \sim\ }
  \operatorname{Hom}_{\operatorname{Alp}/k}
    \bigl(\widehat S_k(E),C\bigr)
\\
  &=
  \operatorname{Hom}_{\operatorname{Vaf}/k}
    \bigl(\operatorname{Spf}(C),\mathbf V^f_k(E)\bigr).
\end{aligned}
\]
In what follows, we identify the formal variety
\(\mathbf V^f_k(E)\) with the \(k\)-functor \(\mathbf V^f_k(E)\).

\medskip
\noindent\textbf{1.2.5.}
\phantomsection
\label{sga3:VIIB:1.2.5}
By \SGARefLink{sga3:vii:paragraph:1-2-4}{1.2.4}, the zero morphism from \(E\) to \(k\) corresponds to a
morphism of profinite algebras
\[
  \pi:\widehat S_k(E)\longrightarrow k.
\]
The morphism \(\pi\) induces the zero map on \(S_k^n(E)\) for
\(n\geq1\), and defines a section of the structure morphism
\[
  \mathbf V^f_k(E)\longrightarrow\operatorname{Spf}(k).
\]

We denote by \(\mathbf V^{f,0}_k(E)\) the formal variety whose points
are the images of the points of \(\operatorname{Spf}(k)\) under the
section \(\operatorname{Spf}(\pi)\), and whose local algebras at
those points are the same as those of \(\mathbf V^f_k(E)\).
\footnote{Editors' note: The remainder of this paragraph has been
spelled out in detail.}

The affine algebra of \(\mathbf V^{f,0}_k(E)\) is thus the separated
completion of \(S_k(E)\) for the topology defined by the ideals
\(\mathfrak l\in\Upsilon\) \textup{(cf.\ \SGARefLink{sga3:vii:paragraph:1-2-4}{1.2.4})} which contain
\(S_k^n(E)\) for all sufficiently large \(n\). Hence it is the
infinite product
\[
  k[[E]]
  =
  k\times E\times S_k^2(E)\times S_k^3(E)\times\cdots.
\]

Let \(C\) be an object of \(\operatorname{Alf}/k\), let
\[
  u\in\mathbf V^f_k(E)(C)=\operatorname{Hom}_c(E,C),
\]
and let
\[
  \phi:\widehat S_k(E)\longrightarrow C
\]
be the corresponding morphism of profinite \(k\)-algebras. Then
\(\operatorname{Ker}(\phi)\) is an open ideal---that is, \(\phi\)
belongs to \(\mathbf V^{f,0}_k(E)(C)\)---if and only if
\(\operatorname{Ker}(\phi)\) contains \(S_k^n(E)\) for all
sufficiently large \(n\), or equivalently if and only if \(u(E)\) is
contained in the radical \(r(C)\) of \(C\). Consequently, for every
object \(C\) of \(\operatorname{Alf}/k\), there are canonical
isomorphisms
\[
  \mathbf V^{f,0}_k(E)(C)
  \simeq
  \operatorname{Hom}_{\operatorname{Alp}/k}(k[[E]],C)
  \simeq
  \operatorname{Hom}_c(E,r(C)).
\]

\medskip
\noindent\textbf{1.2.6.}
\phantomsection
\label{sga3:VIIB:1.2.6}
A formal \(k\)-variety \(V\) is said to be \emph{of finite length} if
its affine algebra has finite length. Similarly, if \(S\) is a scheme,
an \(S\)-scheme \(X\) is said to be \emph{of finite length} if \(X\)
is finite over \(S\) and the direct image of \(\mathcal O_X\) on
\(S\) is an \(\mathcal O_S\)-module of finite length.\footnote{Editors'
note: The following sentence has been added.}
Thus giving an \(S\)-scheme \(X\) of finite length is the same as
giving a finite set
\[
  \{s_1,\ldots,s_n\}
\]
of closed points of \(S\), together with an
\(\mathcal O_{S,s_i}\)-algebra of finite length at each such point.

It follows that the \(S\)-schemes of finite length identify with the
formal varieties of finite length over the following formal scheme
\(\widehat S\). The underlying topological space of \(\widehat S\)
is the set of closed points of \(S\), endowed with the discrete
topology. If \(s\) is such a closed point, the stalk
\(\mathcal O_{\widehat S,s}\) of the structure sheaf is the separated
completion
\[
  \widehat{\mathcal O}_{S,s}
\]
of \(\mathcal O_{S,s}\) for the linear topology defined by the ideals
of finite colength. Thus
\[
  \widehat S
  =
  \operatorname{Spf}\mathcal A(\widehat S),
\]
where \(\mathcal A(\widehat S)\) is the product of the
\(\widehat{\mathcal O}_{S,s}\), as \(s\) ranges over the closed points
of \(S\), endowed with the product topology.

\medskip
\noindent\textbf{Definition.}
\phantomsection
\label{sga3:VIIB:1.2.6.definition}
\textit{If \(X\) is an \(S\)-scheme, let
\(\widehat X/\widehat S\) denote the formal variety over
\[
  k=\mathcal A(\widehat S)
\]
defined as follows. Its underlying topological space consists of the
points \(x\in X\) such that
\[
  [\kappa(x):\kappa(s)]<\infty,
\]
where \(s\) is the image of \(x\) in \(S\). The local ring
\(\mathcal O_{\widehat X/\widehat S,x}\) is the separated completion
of \(\mathcal O_{X,x}\) for the linear topology defined by the ideals
\(I\) of \(\mathcal O_{X,x}\) for which
\(\mathcal O_{X,x}/I\) has finite length as an
\(\mathcal O_{S,s}\)-module. Since
\([\kappa(x):\kappa(s)]<\infty\), this is equivalent to requiring
that \(\mathcal O_{X,x}/I\) have finite length as an
\(\mathcal O_{X,x}\)-module.}

Let \(\operatorname{Vaf}^{\ell f}/\widehat S\) be the category of
formal varieties of finite length over \(\widehat S\), identified
with the category of \(S\)-schemes of finite length. By \SGARefLink{sga3:VIIB:1.1}{1.1} and
\SGARefLink{sga3:VIIB:1.2.1}{1.2.1}, the category \(\operatorname{Vaf}/\widehat S\) of formal
varieties over \(\widehat S\) is equivalent to the category of
left-exact contravariant functors
\[
  \operatorname{Vaf}^{\ell f}/\widehat S
  \longrightarrow
  (\mathbf{Sets}).
\]
In particular, for every \(S\)-scheme \(X\), the functor
\[
  T\longmapsto
  \operatorname{Hom}_{(\operatorname{Sch}/S)}(T,X)
\]
from \(\operatorname{Vaf}^{\ell f}/\widehat S\) to
\((\mathbf{Sets})\) is left exact and therefore corresponds to a
formal variety over \(\widehat S\). This variety is precisely
\(\widehat X/\widehat S\).\footnote{Editors' note: The following
proposition has been expanded by adding the fact that the functor
\(X\mapsto\widehat X/\widehat S\) commutes with finite projective
limits. In the original this appeared in the proof of \SGARefLink{sga3:VIIB:1.3.4}{1.3.4}; the
proof given here is more direct than the original one. This result
will also be used in \SGARefLink{sga3:VIIB:2}{Section~2} and in \SGARefLink{sga3:VIIA:3.3.1}{3.3.1}.}

\medskip
\noindent\textbf{Proposition.}
\phantomsection
\label{sga3:VIIB:1.2.6.proposition}
\textit{For every \(S\)-scheme \(X\), the functors}
\[
  \operatorname{Hom}_{\operatorname{Vaf}/\widehat S}
    (-,\widehat X/\widehat S)
  \qquad\textit{and}\qquad
  \operatorname{Hom}_{(\operatorname{Sch}/S)}(-,X)
\]
\textit{have the same restriction to
\(\operatorname{Vaf}^{\ell f}/\widehat S\). This defines a functor}
\[
  X\longmapsto\widehat X/\widehat S
\]
\textit{from \((\operatorname{Sch}/S)\) to
\(\operatorname{Vaf}/\widehat S\), and this functor commutes with
finite projective limits.}

\noindent\textit{Proof.}
It is immediate that the formal variety \(\widehat X/\widehat S\)
has the required property and that the correspondence
\(X\mapsto\widehat X/\widehat S\) is functorial. It remains to prove
the second assertion.

Put
\[
  \mathscr S=(\operatorname{Sch}/S),
  \qquad
  \mathscr V=\operatorname{Vaf}/\widehat S,
\]
and write \(\widehat X\) for \(\widehat X/\widehat S\).
By \SGARefLink{sga3:VIIB:1.2.B}{1.2.B}, the category \(\mathscr V\) has arbitrary projective
limits. Let \((X_i)_{i\in I}\) be a projective system in
\(\mathscr S\), and suppose that
\[
  X=\varprojlim_iX_i
\]
exists in \(\mathscr S\), as it does when \(I\) is finite. The functor
which associates with every object \(Y\) of \(\mathscr S\) the
functor
\[
  h_Y=\operatorname{Hom}_{\mathscr S}(-,Y)
\]
and, respectively, with every object \(V\) of \(\mathscr V\) the
functor
\[
  h_V=\operatorname{Hom}_{\mathscr V}(-,V)
\]
commutes with projective limits. Hence, for every \(S\)-scheme \(T\)
of finite length, there are functorial isomorphisms
\[
\begin{aligned}
  \operatorname{Hom}_{\mathscr S}(T,X)
  &\simeq
  \varprojlim_i\operatorname{Hom}_{\mathscr S}(T,X_i)
\\
  &\simeq
  \varprojlim_i\operatorname{Hom}_{\mathscr V}(T,\widehat X_i)
  \simeq
  \operatorname{Hom}_{\mathscr V}
    \left(T,\varprojlim_i\widehat X_i\right).
\end{aligned}
\]
Consequently, the functor \(X\mapsto\widehat X/\widehat S\) commutes
with projective limits whenever they exist in
\((\operatorname{Sch}/S)\), and in particular with finite
projective limits.
\hfill\(\square\)

